\def\tituloconcepto{Interferencia cuántica}
% Si el capitulo se compila solo (sin main.tex), esta cabecera abre el documento.
% Antes de % Si el capitulo se compila solo (sin main.tex), esta cabecera abre el documento.
% Antes de % Si el capitulo se compila solo (sin main.tex), esta cabecera abre el documento.
% Antes de \input{cabecera} cada capitulo define \tituloconcepto.
\ifdefined\enmain\else
  \documentclass[11pt,a4paper]{article}
  \input{preambulo}
  \begin{document}
  \title{\tituloconcepto}
  \author{\autor}
  \date{\today}
  \maketitle
\fi
 cada capitulo define \tituloconcepto.
\ifdefined\enmain\else
  \documentclass[11pt,a4paper]{article}
  % Preambulo compartido por todos los documentos. Compilar con XeLaTeX/LuaLaTeX
% (o Tectonic). Si se usa Overleaf, elegir el compilador XeLaTeX.
\usepackage[spanish,es-noshorthands]{babel}
\usepackage{amsmath,amssymb}
\usepackage{booktabs,array}
\usepackage[a4paper,margin=2.5cm]{geometry}
\usepackage{xcolor}
\usepackage{listings}
\usepackage[most]{tcolorbox}
\usepackage{quantikz}
\usepackage[hidelinks]{hyperref}

% ---- Autor (editar a gusto) -------------------------------------------------
\newcommand{\autor}{Genesis Silesky}

% ---- Notacion de Dirac y atajos ---------------------------------------------
% quantikz ya define \ket y \bra; DeclareRobustCommand los redefine sin error.
\DeclareRobustCommand{\ket}[1]{\left|#1\right\rangle}
\DeclareRobustCommand{\bra}[1]{\left\langle#1\right|}
\DeclareRobustCommand{\braket}[2]{\left\langle#1\middle|#2\right\rangle}
\newcommand{\C}{\mathbb{C}}
\newcommand{\R}{\mathbb{R}}
\DeclareMathOperator{\Tr}{Tr}

% ---- Cajas de apoyo --------------------------------------------------------
% "concepto": definiciones previas que hacen falta para entender el capitulo.
\newtcolorbox{concepto}[1]{breakable,colback=blue!4,colframe=blue!55!black,
  fonttitle=\bfseries,title={Concepto previo: #1}}
% "nota": aclaraciones y advertencias.
\newtcolorbox{nota}[1]{breakable,colback=orange!6,colframe=orange!70!black,
  fonttitle=\bfseries,title={#1}}

% ---- Estilos de listados ----------------------------------------------------
\lstdefinestyle{py}{language=Python,basicstyle=\ttfamily\footnotesize,
  breaklines=true,postbreak=\mbox{\textcolor{gray}{$\hookrightarrow$}\space},
  numbers=left,numberstyle=\tiny\color{gray},frame=single,
  commentstyle=\color{green!40!black}\itshape,
  keywordstyle=\color{blue!70!black}\bfseries,
  stringstyle=\color{red!60!black},showstringspaces=false,tabsize=4,
  columns=fullflexible,keepspaces=true}
\lstdefinestyle{salida}{basicstyle=\ttfamily\footnotesize,breaklines=true,
  frame=leftline,backgroundcolor=\color{gray!8},columns=fullflexible,
  keepspaces=true}

  \begin{document}
  \title{\tituloconcepto}
  \author{\autor}
  \date{\today}
  \maketitle
\fi
 cada capitulo define \tituloconcepto.
\ifdefined\enmain\else
  \documentclass[11pt,a4paper]{article}
  % Preambulo compartido por todos los documentos. Compilar con XeLaTeX/LuaLaTeX
% (o Tectonic). Si se usa Overleaf, elegir el compilador XeLaTeX.
\usepackage[spanish,es-noshorthands]{babel}
\usepackage{amsmath,amssymb}
\usepackage{booktabs,array}
\usepackage[a4paper,margin=2.5cm]{geometry}
\usepackage{xcolor}
\usepackage{listings}
\usepackage[most]{tcolorbox}
\usepackage{quantikz}
\usepackage[hidelinks]{hyperref}

% ---- Autor (editar a gusto) -------------------------------------------------
\newcommand{\autor}{Genesis Silesky}

% ---- Notacion de Dirac y atajos ---------------------------------------------
% quantikz ya define \ket y \bra; DeclareRobustCommand los redefine sin error.
\DeclareRobustCommand{\ket}[1]{\left|#1\right\rangle}
\DeclareRobustCommand{\bra}[1]{\left\langle#1\right|}
\DeclareRobustCommand{\braket}[2]{\left\langle#1\middle|#2\right\rangle}
\newcommand{\C}{\mathbb{C}}
\newcommand{\R}{\mathbb{R}}
\DeclareMathOperator{\Tr}{Tr}

% ---- Cajas de apoyo --------------------------------------------------------
% "concepto": definiciones previas que hacen falta para entender el capitulo.
\newtcolorbox{concepto}[1]{breakable,colback=blue!4,colframe=blue!55!black,
  fonttitle=\bfseries,title={Concepto previo: #1}}
% "nota": aclaraciones y advertencias.
\newtcolorbox{nota}[1]{breakable,colback=orange!6,colframe=orange!70!black,
  fonttitle=\bfseries,title={#1}}

% ---- Estilos de listados ----------------------------------------------------
\lstdefinestyle{py}{language=Python,basicstyle=\ttfamily\footnotesize,
  breaklines=true,postbreak=\mbox{\textcolor{gray}{$\hookrightarrow$}\space},
  numbers=left,numberstyle=\tiny\color{gray},frame=single,
  commentstyle=\color{green!40!black}\itshape,
  keywordstyle=\color{blue!70!black}\bfseries,
  stringstyle=\color{red!60!black},showstringspaces=false,tabsize=4,
  columns=fullflexible,keepspaces=true}
\lstdefinestyle{salida}{basicstyle=\ttfamily\footnotesize,breaklines=true,
  frame=leftline,backgroundcolor=\color{gray!8},columns=fullflexible,
  keepspaces=true}

  \begin{document}
  \title{\tituloconcepto}
  \author{\autor}
  \date{\today}
  \maketitle
\fi


\section{Interferencia cuántica}

\subsection{Objetivo}
Mostrar que las amplitudes se suman y pueden cancelarse, y que por eso el
resultado de un circuito depende de las \emph{fases} y no solo de las
probabilidades intermedias.

\subsection{Conceptos previos}
\begin{concepto}{amplitudes con signo y suma de amplitudes}
Ver sección~\ref{sec:prerreq}. Una amplitud puede ser negativa o compleja;
una probabilidad no. Al encadenar puertas, las amplitudes que llegan al mismo
resultado se \emph{suman} (con su signo/fase) \emph{antes} de elevar al
cuadrado.
\end{concepto}

\subsection{Fundamento teórico}
Aplicando $H$ dos veces a $\ket0$:
\begin{align}
  H\ket0&=\tfrac1{\sqrt2}(\ket0+\ket1),\\
  HH\ket0&=\tfrac1{\sqrt2}\Big[\tfrac1{\sqrt2}(\ket0+\ket1)+
     \tfrac1{\sqrt2}(\ket0-\ket1)\Big]=\ket0 .
\end{align}
La amplitud de $\ket1$ es $\tfrac12-\tfrac12=0$ (interferencia
\emph{destructiva}); la de $\ket0$ es $\tfrac12+\tfrac12=1$ (\emph{constructiva}).
Un modelo clásico probabilístico (50/50 y luego 50/50) daría de nuevo
50/50, porque las probabilidades nunca se cancelan; las amplitudes sí.

Con una fase en medio, $Z$ cambia el signo de $\ket1$:
$ZH\ket0=\tfrac1{\sqrt2}(\ket0-\ket1)=\ket-$ y $H\ket-=\ket1$. En general
$HZH=X$.
\begin{center}
\begin{quantikz}
\lstick{$\ket0$} & \gate{H} & \gate{Z} & \gate{H} & \meter{}
\end{quantikz}
\end{center}
Este es el principio de los algoritmos cuánticos: diseñar fases para que las
respuestas incorrectas se cancelen y la correcta se refuerce (capítulo~06).

\subsection{Código (completamente comentado)}
\lstinputlisting[style=py]{../../fuentes/05a_interferencia.py}

\subsection{Explicación celda por celda}
\begin{description}
  \item[Celda 2] Compara un circuito con una $H$ (50/50) y otro con dos $H$
    (solo $\ket0$). Se usa \texttt{Statevector} para ver las amplitudes
    exactas, sin aleatoriedad.
  \item[Celda 3] Inserta $Z$ entre las dos $H$ y mide 1000 shots: sale
    siempre ``1''. Solo se cambió una fase.
\end{description}

\subsection{Salida obtenida}
\lstinputlisting[style=salida]{../../fuentes/salidas/05a_interferencia.txt}
Es determinista porque la medición final tiene probabilidad $1$.

\subsection{Interpretación}
La puerta $Z$ ``no hace nada'' a las probabilidades de $\ket+$, pero cambia
por completo el resultado final. Eso solo es posible porque la información
de fase se almacena en el signo de las amplitudes.

\subsection{Ejercicios}
\begin{enumerate}
  \item Reemplace $Z$ por la puerta $S$ (fase $i$) y prediga el resultado
    antes de ejecutarlo.
  \item Pruebe $H\,X\,H$: ¿a qué puerta equivale?
  \item Explique por qué un ``coin flip'' clásico repetido no reproduce el
    resultado de $HH$.
\end{enumerate}

\ifdefined\enmain\else
  \end{document}
\fi

